\documentclass[]{article}

%opening
\title{Add a screen to your micro controller with an IDEC HMI!}
\author{Senju}

\begin{document}
	
	\maketitle
	
	\begin{abstract}
		There are many Micro Controller development boards capable of handling video, yet many popular models have no such capability. Arduino UNO, Mega, teensy, to name a few.
		IDEC HMIs have a unique approach to its programming which makes it possible to add a screen to projects which ordinarily would be unable to support one. 
	\end{abstract}
	\section*{Motivation} This project could be achieved in a variety of ways, using the same hardware, with different software configurations, or by utilizing different hardware. 
	I happen to have several IDEC HG2J HMIs laying around, so I used them. 
	In this implementation, we utilize Maintenance Protocol the native communication for IDEC products. Adding a screen to an arduino Mega, using maintenance protocol on the HMI means we 
	can skip the step of pawing through HMI configurations, and get right to knocking out a project.  I coded this in C, but other languages could work.  
	
	\section*{How does it work?}
	The hack of using the HMI is that the HMI uses a preformed GUI to populate screen elements. The Arduino doesn't have to handle the visual components at all, it just sends data to the HMI to be rendered.
	Because of how the IDEC HMI works, independent of data stream, the HMI reads registers, and displays them on screen. 
	\section*{implementation}
	To build this project, or one like it you will need:
	An IDEC HMI
	An Arduino
	A network connection, I used a WIZ1000 Ethernet shield.
	An ethernet cable. 
	
	Nice to have: a Ethernet switch, additional ethernet cables. 
	
	\section*{set up}
	I would call this initialization, but IDEC uses the term initialization to mean factory reset.
	Plug 24 VDC into the HMI.
	Design your GUI.
	Download your GUI.
	
	\section*{supporting software}
	I used Geany to write the code base, with Arduino build and compile macros. 
	NV4 was used to build and download the GUI
	\section*{other methods}
	Other methods using the same hardware exist; you could implement similar functionality using modbus, mqtt,  FTP or serial, but I find the other methods of implementation tedious. 
	\section*{Notes:}
	Current implementation is read only. The HMI requests data from the Arduino.
	While it is possible to send instructions, to the Arduino, I haven't gotten that far yet.
	
	\section*{}
	\section*{}
	
\end{document}
